\documentclass[10pt,a4paper]{amsart}
\usepackage[margin=19mm]{geometry}
\usepackage{amsmath,amssymb,mathtools}
\usepackage[scr=rsfs,cal=euler]{mathalfa}
\usepackage{xfrac}
\usepackage[unicode,hidelinks,bookmarks=false]{hyperref}
\newcommand{\ie}{i.e.\ }
\newcommand{\cf}{cf.\ }
\newcommand{\resp}{respectively\ }
\newcommand{\Subsection}{\subsection}
\providecommand{\nobreakdash}{\nobreak\hbox{-}}
\setlength{\emergencystretch}{2em}
\multlinegap0pt
\usepackage{amssymb}
\newtheorem{theorem}{Theorem}
\title{A Simple Trigonometric Classification of Quartic Roots}
\author{Sawon Pratiher}
\date{Source: arXiv:2603.28814v1 (CC BY 4.0)}
\newtheorem{lemma}{Lemma}
\newtheorem{corollary}{Corollary}
\newtheorem{proposition}{Proposition}
\begin{document}
\maketitle
\par\medskip\noindent\textbf{Source applicability on hold.} The recorded source range or claim issue is not resolved by this context repair. The author text is retained; consult the canonical metadata for the precise issue and evidence.\par\medskip

\par\medskip\noindent\textbf{Current collection classification: H1.} Retained author material; this item does not qualify for the advanced H2/H3 collection. See the canonical metadata for the source-specific rationale.\par\medskip


{\sloppy \noindent\emph{Source and licence.} A Simple Trigonometric Classification of Quartic Roots. Version arXiv:2603.28814v1 (CC BY 4.0), distributed by arXiv under the Creative Commons Attribution 4.0 International licence, \url{http://creativecommons.org/licenses/by/4.0/}, as recorded in the arXiv licence metadata for this exact version and shown on the abstract page \url{https://arxiv.org/abs/2603.28814v1.} Full licence notice: \path{licenses/arxiv-2603.28814-CC-BY-4.0.txt}.\par\smallskip\par}

\noindent\textit{Editorial scope.} Counted unit: the article's theorem thm:four\_complex (source0.tex lines 165--170). The article's complete original proof of it is delivered with it (source0.tex lines 172--178, one \texttt{proof} environment of 7 lines). Additional author definitions and conventions recorded in the correction metadata are retained below as uncounted context. Source caveats and pending qualification issues are recorded separately. Author statements and proofs remain unchanged; the fixed complete author source remains bundled.\par\medskip
\par\medskip\noindent\textit{Retained author context (same fixed source; added for the selected problem).}\par
\begin{equation}\label{eq:cheb4}
8\cos^4\theta - 8\cos^2\theta + 1 = \cos 4\theta,
\end{equation}

\begin{equation}\label{eq:depressed}
P(t) = t^4 + mt^2 + pt + q = 0.
\end{equation}

Here is the key step. Assume $m < 0$, and set $u = \sqrt{-m}$. Substituting $t = u\cos\theta$ into (\ref{eq:depressed}) and dividing through by $u^4/8 = m^2/8$ gives
\begin{equation}\label{eq:after_sub}
8\cos^4\theta - 8\cos^2\theta + \frac{8p}{(-m)^{3/2}}\cos\theta + \frac{8q}{m^2} = 0.
\end{equation}
Observe that the first two terms are exactly the left side of the Chebyshev identity (\ref{eq:cheb4}). Subtracting that identity from (\ref{eq:after_sub}), we arrive at the \textbf{reduced trigonometric equation}:
\begin{equation}\label{eq:f_def}
\boxed{f(\theta) \;=\; a\cos\theta + \cos 4\theta + b \;=\; 0,}
\end{equation}
where the \textbf{trigonometric parameters} are
\begin{equation}\label{eq:ab}
a = \frac{8p}{(-m)^{3/2}}, \qquad b = \frac{8q}{m^2} - 1.
\end{equation}

Since $\cos\colon [0,\pi] \to [-1,1]$ is a continuous, strictly decreasing bijection, the map $\theta \mapsto u\cos\theta$ is a bijection from $[0,\pi]$ onto $[-u, u]$. A direct calculation confirms the identity
\begin{equation}\label{eq:fP}
f(\theta) = \frac{8\,P(u\cos\theta)}{u^4}
\end{equation}
for all $\theta$, so the zeros of $f$ on $[0,\pi]$ correspond exactly to the roots of $P$ in $[-u, u]$. In particular, evaluating at the endpoints:

\begin{lemma}[{Convexity outside $[-u,u]$}]\label{lem:convex}
If $m < 0$, then $P$ is strictly convex on $(-\infty, -u)$ and on $(u, \infty)$.
\end{lemma}

\begin{proof}
The second derivative is $P''(t) = 12t^2 + 2m$. For $|t| > u = \sqrt{-m}$, we have $t^2 > -m$, so $P''(t) > 12(-m) + 2m = -10m > 0$.
\end{proof}

A strictly convex function can cross the horizontal axis at most twice, but since $P(t) \to +\infty$ as $t \to \pm\infty$ for a quartic with positive leading coefficient, it actually crosses at most \emph{once} on each half-line. This yields:

\begin{corollary}\label{cor:ext}
The quartic $P$ has at most one root in $(u, \infty)$ and at most one root in $(-\infty, -u)$. Moreover:
\begin{itemize}
\item $P$ has exactly one root in $(u, \infty)$ if and only if $P(u) < 0$, i.e., $f(0) < 0$.
\item $P$ has exactly one root in $(-\infty, -u)$ if and only if $P(-u) < 0$, i.e., $f(\pi) < 0$.
\item $P$ has no roots outside $[-u,u]$ if and only if $f(0) \geq 0$ and $f(\pi) \geq 0$.
\end{itemize}
\end{corollary}

\begin{proof}
Since $P(t) \to +\infty$ as $t \to +\infty$ and $P$ is strictly convex and continuous on $(u, \infty)$, the intermediate value theorem guarantees exactly one crossing if $P(u) < 0$, and no crossing if $P(u) \geq 0$. The argument for $(-\infty, -u)$ is identical.
\end{proof}

We are now ready to state the main results. Define
\[
N_{\mathrm{int}} = \#\{\theta \in [0,\pi] : f(\theta) = 0\}, \qquad N_{\mathrm{ext}} = \mathbf{1}_{f(0)<0} + \mathbf{1}_{f(\pi)<0},
\]
where $\mathbf{1}$ denotes the indicator function. By Proposition~\ref{prop:bijection_formal} and Corollary~\ref{cor:ext}, the total number of real roots of $P$ is
\begin{equation}\label{eq:Nreal}
N_{\mathrm{real}} = N_{\mathrm{int}} + N_{\mathrm{ext}}.
\end{equation}

\begin{proposition}\label{prop:bijection_formal}
The zeros of $f$ in $[0,\pi]$ are in bijection with the roots of $P$ in $[-u,u]$. There is no double-counting at the boundary: if $f(0) = 0$ then $t = u$ is counted in $N_{\mathrm{int}}$, and $N_{\mathrm{ext}}$ contributes zero from the right (and similarly for $f(\pi) = 0$).
\end{proposition}

\begin{equation}\label{eq:boundary}
f(0) = \frac{8\,P(u)}{u^4}, \qquad f(\pi) = \frac{8\,P(-u)}{u^4}.
\end{equation}

\begin{theorem}[All four roots complex]\label{thm:four_complex}
Suppose $m < 0$. All four roots of $P(t) = 0$ are complex if and only if $f(\theta) > 0$ for all $\theta \in [0,\pi]$ (equivalently, $N_{\mathrm{int}} = 0$ and $N_{\mathrm{ext}} = 0$). A sufficient condition is
\begin{equation}\label{eq:four_complex_suff}
b > |a| + 1.
\end{equation}
\end{theorem}

\begin{proof}
If $f > 0$ on $[0,\pi]$, then $N_{\mathrm{int}} = 0$. The boundary values $f(0)$ and $f(\pi)$ are both positive, so $P(u) > 0$ and $P(-u) > 0$ by (\ref{eq:boundary}), giving $N_{\mathrm{ext}} = 0$. Thus $N_{\mathrm{real}} = 0$.

For the sufficient condition: for all $\theta$, $|a\cos\theta + \cos 4\theta| \leq |a| + 1$, so $f(\theta) \geq b - (|a|+1) > 0$ when $b > |a| + 1$.

Conversely, if $f < 0$ on $[0,\pi]$, then $N_{\mathrm{int}} = 0$ but $f(0) < 0$ and $f(\pi) < 0$, giving $N_{\mathrm{ext}} = 2$, so there are two (not zero) real roots. Thus the only route to zero real roots is $f > 0$.
\end{proof}

\end{document}
